%表紙
\newpage
\chapter{解任・退会規則}
\begin{enumerate}[label=第\arabic*条　]
  \item この章は、音楽ゲーム学園学園憲章（以下「憲章」という。）の趣旨に基づき、役職者の辞任及び解任並びに学生（役職者を含む。）に対する処分に関し必要な事項を定めるものとする。
  \item 学園長を除く役職者は、学園長への届出により、いつでもその職を辞することができる。
    \begin{enumerate}[label=\arabic*　,start=2,leftmargin=*]
      \item 学園長は、教務主事への通知により、いつでもその職を辞することができる。
    \end{enumerate}
  \item 学園長は、学園長を除く役職者が次の各号のいずれかに該当すると認めるときは、教務主事の協議を経て、当該役職者をその職から解くことができる。
      \begin{enumerate}[label=(\arabic*)　]
        \item その職務を行うことができず、学園の円滑な運営に支障を生じているとき
        \item 他の学生に対するハラスメントその他の行為により、役職者としての職務を託するに適しないと認められるとき
      \end{enumerate}
    \begin{enumerate}[label=\arabic*　,start=2,leftmargin=*]
      \item 学園長が前項各号のいずれかに該当すると認めるときは、教務主事は、その協議により、当該学園長をその職から解くことができる。
      \item 前二項の規定による解任は、この章に定める処分に代わるものではなく、学園長又は教務主事は、必要があると認めるときは、別途この章の定めるところにより処分を行うことができる。
    \end{enumerate}
  \item 学園長又は教務主事は、次条各号のいずれかに該当する事実があると認めるときは、当該学生に対し、次の各号に掲げる処分を行うことができる。
      \begin{enumerate}[label=(\arabic*)　]
        \item 警告
        \item 活動停止（期間を定めて、学園における活動の全部又は一部を停止させることをいう。）
        \item 除名
      \end{enumerate}
  \item 学生が次の各号のいずれかに該当する行為を行ったときは、前条の処分の対象とすることができる。
      \begin{enumerate}[label=(\arabic*)　]
        \item 憲章又は学園規則に違反する行為
        \item 他の学生に対するハラスメントその他の著しく不当な言動
        \item 各種の申請又は届出における虚偽の記載その他の不正な行為
        \item 学園の名誉又は信用を著しく毀損する行為
        \item 学園が活動の基盤として利用する外部サービスの利用条件に違反する行為であって、学園の活動に支障を生じさせるもの
        \item 前各号に掲げるもののほか、他の学生の権利利益を著しく損なう行為
      \end{enumerate}
  \item 処分は、警告、活動停止、除名の順に、段階的に行うことを原則とする。
    \begin{enumerate}[label=\arabic*　,start=2,leftmargin=*]
      \item 学園長又は教務主事は、警告を受けた学生が当該学期内に再び前条各号のいずれかに該当する行為を行ったと認めるときは、活動停止の処分を行うことができる。
      \item 学園長又は教務主事は、活動停止の処分を受けた学生が、当該処分の期間中又はその終了後に再び前条各号のいずれかに該当する行為を行ったと認めるときは、除名の処分を行うことができる。
    \end{enumerate}
  \item 前条の規定にかかわらず、著しく重大な違反があると認めるときは、学園長又は教務主事は、警告を経ることなく、直ちに活動停止又は除名の処分を行うことができる。
  \item 学園長又は教務主事は、学生の行為により学園の秩序又は他の学生の安全が現に害され、又は害されるおそれが急迫していると認めるときは、前各条の処分によらず、必要な範囲で当該学生の学園における活動へのアクセスを一時的に制限する等の緊急の措置を講ずることができる。
    \begin{enumerate}[label=\arabic*　,start=2,leftmargin=*]
      \item 前項の措置を講じたときは、学園長又は教務主事は、遅滞なく、当該学生に対しその理由を通知しなければならない。
      \item 第１項の措置は、当該措置を講じた日の属する学期の末日までにその効力を失う。
      \item 第１項の措置は、この章に定める処分に代わるものではなく、学園長又は教務主事は、必要があると認めるときは、別途この章の定めるところにより処分を行うことができる。
    \end{enumerate}
  \item 学生は、第３条の規定による解任、前条の措置又はこの章に定める処分について、学園長又は教務主事に対し、不服を申し立てることができる。
    \begin{enumerate}[label=\arabic*　,start=2,leftmargin=*]
      \item 前項の不服申立てについては、当該解任、措置又は処分を行った者以外の役職者が、その審査に当たるものとする。
    \end{enumerate}
  \item この章の実施に関し必要な事項は、学園長が別に定める。
\end{enumerate}